\documentclass[12p, oneside]{book}
\usepackage[T1]{fontenc}
\usepackage[polish]{babel}
\usepackage{color}
\usepackage[table, usenames, dvipsnames]{xcolor}
\usepackage{hyperref}
\usepackage{graphicx}
\hypersetup{
    colorlinks=true, 
    linktoc=all,    
    linkcolor=blue,
}



\newcommand{\steno}[1]{\textbf{\textsf{\textcolor{violet}{ #1}}}}

\newcommand{\stenoExc}[1]{\textbf{\textsf{\textcolor{red}{ #1}}}}

\newenvironment{tabela2}
	{
		\rowcolors{1}{cyan}{Emerald}
		\begin{tabular}{c c}
	}
	{
		\end{tabular}
	}


\setcounter{secnumdepth}{0}
\begin{document}

\frontmatter

\title{Polski System Stenotypii}

\author{B. P.}
\maketitle
\newpage

\tableofcontents

\mainmatter

\chapter{Wstęp}

\chapter{Kurs}

\chapter{Pełny opis systemu}
\section{Klawiatura}
\includegraphics[width=12cm]{grafika/pss/klawiatura.jpg}

Powyżej przedstawiony jest układ klawiatury wykorzystywany przez PSS. Klawiatura jest podzielona na cztery części:
\begin{itemize}
	\item Lewa górna --- odpowiedzialna za spółgłoski początkowe oraz samogłoski posiłkowe --- tę część obsługuje się palcami lewej dłoni
	\item Lewa dolna --- klawisz J odpowiedzialny za udźwięcznianie oraz dwa klawisze samogłosek --- tę część obsługuje się lewym kciukiem
	\item Prawa dolna --- trzy klawisze samogłosek --- tę część obsługuje się prawym kciukiem
	\item Prawa górna --- odpowiedzialna za spółgłoski końcowe oraz samogłoski dodatkowe --- tę część obsługuje się palcami prawej dłoni
\end{itemize}

\includegraphics[width=12cm]{grafika/pss/pozycjepoczatkowe.jpg}
Początkowe pozycje poszczególnych palców razem z obsługiwanymi przez nie klawiszami.

\section{Budowa akordu}
Akord --- jednoczesne wciśnięcie i puszczenie pewnej ilości klawiszy na klawiaturze w taki sposób, aby wszystkie zostały zarejestrowane i przetworzone przez oprogramowanie stenograficzne.
\subsection{Porządek stenotypiczny}
Akordy zapisuje się, wymieniając wykorzystane klawisze w następującym porządku:

\steno{X F Z S K T P V L R J E I \# * Y A U C R L B S G T W O V}

W przypadku niewystąpienia wewnątrz akordu żadnego z klawiszy z części środkowej (tj. \steno{J, E, I, \#, *, Y, A} oraz \steno{U}) w jego miejsce wpisywany jest myślnik (\steno{--}).
\subsection{Rdzenie}
Rdzeń to najbardziej podstawowa forma akordu. Co do zasady składa się z trzech części:
\begin{enumerate}
	\item Spółgłoski początkowe
	\item Samogłoski
	\item Spółgłoski końcowe
\end{enumerate}
Oprócz tego w wielu przypadkach (w celu zmniejszenia liczby akordów) będą wykorzystywane \hyperref[sec:dodatkowe]{samogłoski dodatkowe}, których zasady wykorzystywania są przedstawione w rozdziale im poświęconym.

Rdzenie najczęściej pokrywają się z semantycznymi rdzeniami wyrazów. 

Domyślnie rdzenie nie łączą się ze sobą.

\subsubsection{Samogłoski}
Samogłoski stanowią ośrodek prawie każdego rdzenia. Są zapisywane z wykorzystaniem kciuków. Samogłoski przyjmują dwie postacie --- podstawową i zmiękczającą --- w zależności od obecności klawisza \steno{J} w grupie spółgłosek początkowych. Szczegóły na temat zapisywania \hyperref[sec:samogloski]{samogłosek} można znaleźć wewnątrz rozdziału im poświęconego.

\subsubsection{Spółgłoski początkowe}
Bardzo często właśnie od spółgłosek początkowych zależy znaczenie wyrazu. Są to spółgłoski występujące \emph{przed} samogłoskami. Każdej możliwej grupie spółgłoskowej przypisana jest odpowiadająca jej kombinacja klawiszy po lewej stronie klawiatury. Ich listę, jak również szczegółowe zasady wykorzystywania można zobaczyć w rozdziale poświęconym \hyperref[sec:poczatkowe]{spółgłoskom początkowym}.

\subsubsection{Spółgłoski końcowe}
Spółgłoski końcowe są zapisywane palcami prawej dłoni. Są to spółgłoski występujące \emph{po} samogłosce w danym akordzie. Lista poszczególnych kombinacji dostępna jest w rozdziale poświęconym \hyperref[sec:koncowe]{spółgłoskom końcowym}.

\subsection{Przyrostki}
Przyrostki to akordy, których przeznaczeniem jest połączyć się z poprzednim podanym akordem. Są to zwykłe rdzenie, do których dołączono specjalną cząstkę zgodnie z następującą tabelą:

\begin{tabular}{| c | c |}
	\hline
	Klawisze & Od czego zaczyna się przyrostek \\ \hline
	\steno{ZKPL--} & Od samogłoski \\ \hline
	\steno{--OY} & Od klawiszy \steno{Z--} lub \steno{S--} \\ \hline
	\steno{XF--} & Wszystkie pozostałe \\ \hline
\end{tabular}


\subsubsection{Szczególne zachowania}

Ciekawym zachowaniem, właściwym dla przyrostków w PSS jest zachowanie w obecności rdzeniów kończących się na głoskę zmiękczoną (ś, ź, ć, ń, dź). Po przyłączeniu przyrostka zaczynającego się od samogłoski, cała grupa zostanie zmieniona na odpowiadającą jej ze zmiękczeniem przez i niezgłoskotwórcze. Przykłady:

\begin{center}
troć + -ie $\rightarrow$ trocie
		
brać + -i $\rightarrow$ braci

koń + -e $\rightarrow$ konie
\end{center}

Umożliwia to pewną elastyczność co do używania na końcu rdzeni głosek zmiękczonych lub podstawowych.

Ze względów etymologicznych, podobne zjawisko ma miejsce przy rdzeniach kończących się na t i przyrostkach zaczynających się od i:

\begin{center}
brat + -ia $\rightarrow$ bracia

kot + -ie $\rightarrow$ kocie
\end{center}

Wyrazy, których nie dotyczy ta zasada (sympatia, apatia etc.) są uwzględnione w słowniku i nie podlegają temu zjawisku automatycznie.

\subsection{Przedrostki}
Przedrostki to akordy, które mają połączyć się z kolejnym podanym akordem --- odpowiadają semantycznym przedrostkom.

Tworzy się je, dodając do standardowego rdzenia znak \steno{*}.

Należy być ostrożnym z tworzeniem przedrostków, gdyż istotna część z nich jest skróceniami.

\subsection{Specjalne}
Istnieją akordy o specjalnym znaczeniu:

\begin{tabular}{| c | c |}
	\hline
	Akord & Znaczenie \\ \hline
	\stenoExc{XF--OV} & Połączenie dwóch rdzeni bez spacji \\ \hline
	\stenoExc{XF}\hyperref[sec:samogloski]{samogłoska}\stenoExc{OV} & Połączenie dwóch rdzeni bez spacji, ze wstawieniem odpowiedniego wrostka \\ \hline
	\stenoExc{XF*OV} & Wymuszenie spacji \\ \hline
\end{tabular}

\section{Samogłoski}
\label{sec:samogloski}
Przez samogłoskę wewnątrz PSS są rozumiane wszystkie samogłoski, dyftongi, oraz ich połączenia z "i niezgłoskotwórczym". Każda samogłoska ma dwie wersje --- podstawową i udźwięczniającą. Udźwięczniająca jest wykorzystywana w połączeniu ze wszystkimi grupami spółgłosek początkowych zawierających klawisz \steno{J} (nazwa pochodzi od faktu, że w większości przypadków \steno{J} wewnątrz grupy spółgłosek początkowych oznacza udźwięcznienie w stosunku do grupy bez niego).


\section{Spółgłoski początkowe}
\label{sec:poczatkowe}

\section{Spółgłoski końcowe}
\label{sec:koncowe}

\section{Samogłoski dodatkowe}
\label{sec:dodatkowe}

\section{Samogłoski posiłkowe}
\label{sec:posilkowe}
\chapter{Przykłady}
Poniżej dostarczone są przykładowe teksty z proponowanym podziałem na akordym, zapisem stenotypicznym prostym oraz ze skróceniami. Należy pamiętać, że jest to jedynie propozycja, aczkolwiek ściśle spójna z przedstawionym systemem. Pomimo to, w przypadku znalezienia szybszego sposobu zapisu należy się go trzymać.
\section{Przykład 1 --- fragment deklaracji praw człowieka}
\subsection{Tekst}

Zgromadzenie Ogólne ogłasza niniejszą Powszechną Deklarację Praw Człowieka jako wspólny standard do osiągnięcia przez wszystkie ludy i wszystkie narody, tak aby każda jednostka i każdy organ społeczny, mając stale na względzie niniejszą Deklarację, dążyły poprzez nauczanie i wychowanie do zapewnienia poszanowania tych praw i wolności, jak również do zabezpieczenia, poprzez rozwój środków o charakterze krajowym i międzynarodowym, ich powszechnego i skutecznego uznania i przestrzegania, tak pośród ludów Państw Członkowskich, jak też pośród ludów terytoriów znajdujących się pod ich jurysdykcją. 

\subsection{Podział na akordy --- podstawowy}
Zgrom  /  -adz / -en / -ie / O- / gól / -ne / o- / głasz / -a / nin / -iej / -szą / Po- / wszechn / -ą / De- / klar / -ac / -ję / Praw / Człow / -ie / -ka / jak / -o / wspól / -ny / stand / -ard / do / o- / siągn / -ięc / -ia / przez / wszystk / -ie / lud / -y / i / wszystk / -ie / nar / -od / -y / , / tak / ab / -y / każd / -a / jedn / -o- / -stka / i / każd / -y / org / -an / społ / -ecz / -ny / , / maj / -ąc / stal / -e / na / względz / -ie / nin / -iej / -szą / De- / klar / -ac / -ję / , / dąż / -ył / -y / po- / przez / naucz / -an / -ie / i / wy- / chow / -an / -ie / do / za- / pew / -nie / -nia / po- / szan / -ow / -an / -ia / tych / praw / i / woln / -ośc / -i / , / jak / rów / -nież / do / za- / bezp / -iecz / -eni / -a / , / po- / przez / rozw / -ój / środk / -ów / o / char / -a- / kterz / -e / kraj / -ow / -ym / i / międz / -y- / nar / -od / -ow / -ym / , / ich / po- / wszechn / -eg / -o / i / skut / -eczn / -eg / -o / u- / znan / -ia / i / prze- / strzeg / -an / -ia / , / tak / po- / śród / lud / -ów / Państw / Członk / -ow / -skich / , / jak / też / po- / śród / lud / -ów / ter / -yt / -or / -iów / znajd / -uj / -ąc / -ych / się / pod / ich / jur / -y- / sdyk / -cją / . 

\end{document}
